% Figure: a linearly polarised plane wave. E in the x-z plane (engblue),
% B in the y-z plane (engred), both transverse to the propagation axis z,
% in phase, with the same spatial period. The triplet (E, B, k) is right-handed.
\begin{figure}[htbp]
\centering
\begin{tikzpicture}[scale=1.0, >=latex]
  % propagation axis z
  \draw[->, thick] (0,0,0) -- (10,0,0) node[anchor=west] {$z$ (propagation, $\vect{k}$)};
  % E axis (x, vertical)
  \draw[->, thick] (0,0,0) -- (0,2.6,0) node[anchor=south] {$x$ ($\vect{E}$)};
  % B axis (y, into the page diagonal)
  \draw[->, thick] (0,0,0) -- (0,0,3.2) node[anchor=north east] {$y$ ($\vect{B}$)};
  % E field: sine in the x-z plane
  \draw[engblue, thick, samples=80, domain=0:9, variable=\t]
    plot ({\t}, {2.0*sin(\t*40)}, {0});
  % E field arrows (sampled)
  \foreach \t in {0.4,1.2,2.0,2.8,3.6,4.4,5.2,6.0,6.8,7.6,8.4}{
    \draw[engblue, ->] (\t,0,0) -- (\t,{2.0*sin(\t*40)},0);
  }
  % B field: sine in the y-z plane, same phase, smaller drawn amplitude
  \draw[engred, thick, samples=80, domain=0:9, variable=\t]
    plot ({\t}, {0}, {1.6*sin(\t*40)});
  \foreach \t in {0.4,1.2,2.0,2.8,3.6,4.4,5.2,6.0,6.8,7.6,8.4}{
    \draw[engred, ->] (\t,0,0) -- (\t,0,{1.6*sin(\t*40)});
  }
  % wavelength marker (one period spans 9 deg-units => period 360/40 = 9)
  \draw[<->, enggray] (0,-0.5,0) -- (9,-0.5,0)
    node[midway, below, font=\scriptsize] {$\lambda$};
  \node[engblue, font=\scriptsize] at (9.3,1.2,0) {$\vect{E}$};
  \node[engred, font=\scriptsize] at (9.3,0,2.2) {$\vect{B}$};
\end{tikzpicture}
\caption[A linearly polarised plane electromagnetic wave]{A linearly
polarised plane electromagnetic wave travelling along $z$. The electric
field $\vect{E}$ oscillates in the $x$ direction (blue), the magnetic
field $\vect{B}$ oscillates in the $y$ direction (red), the two are in
phase and mutually perpendicular, and both are transverse to the
propagation direction $\vect{k}$. The amplitudes satisfy $|\vect{B}| =
|\vect{E}|/c$; the drawn magnetic amplitude is scaled for legibility.
The triplet $(\vect{E}, \vect{B}, \vect{k})$ forms a right-handed set.}
\label{fig:vol04ch10:plane-wave-eb}
\end{figure}
